\chapter{101}



Diesmal ging sie alleine hin. Sie dachte an Louis. Beim Gedanken an ihn wurde ihr warm.
Seit Giorgias Sturz fühlte sich Valentina ihr wieder nah wie früher.
Jede Rivalität war erloschen.

Sie verliess die Metro-Station über die Treppe nach oben. Die Sonne
schlug ihr mit beispielloser Hitze entgegen. Der Mann war ihr gleich
aufgefallen. Er hatte einige Sitze hinter ihr gesessen. Beim Aussteigen
folgte er ihr. Sie blickte zurück. Dann sah sie es. Er schloss für einen
Wimpernschlag die Augen und nickte ihr zu.
Sie würde sich daran gewöhnen müssen. Und hoffte inbrünstig, das Theater
möge bald zu Ende sein.
Wie zweischneidig sich alles anfühlte. Da lag ihr einerseits die Angst
über die Drohung der beiden Männer im Nacken, andererseits
hatte es wirklich etwas Aufregendes, rund um die Uhr beschützt zu werden. Sie
sollte ihr Denk-Karussell abschalten. In wenigen Minuten würde sie
Giorgia sehen.

Bald lagen die Gebäude des \emph{«San Raffaele»} vor ihr. Beim Eintreten
in die Empfangshalle des \emph{Ospedale} suchte sie nach innerer 
Ruhe. Wie würde sie ihre Freundin antreffen? Wie beim letzten Mal fühlten sich ihre Beine schwer an. Nicht zu wissen, ob sie überhaupt ansprechbar war, war das Schlimmste. Sie erinnerte sich an die Worte des Arztes.

«Kommen Sie, so oft Sie können. Auch wenn Signora Sala Sie nicht bewusst
wahrnimmt, Sie werden ihr guttun. Versuchen Sie ruhig zu bleiben.
Sprechen Sie langsam und erzählen Sie von gemeinsamen, schönen Dingen.»

Valentina hatte ihn angesehen und nichts zu sagen gewusst.

Der lange Korridor lag in einem schweren Duftgemisch von
Desinfektionsmitteln. Kein Mensch weit und breit. Jeder ihrer Schritte
kam ihr vor wie gewaltsames Eindringen in eine verlorene Welt.
Valentina versuchte, leise aufzutreten. Sie krallte ihre Flipflops mit
den Zehen an die Fusssohlen.
Dann klopfte sie kurz und öffnete die Zimmertür. Über dem Raum lag ein
Schleier, gemacht aus einem warmen Gemisch aus Krankheit und
frisch gewaschenen Tüchern. Sie atmete durch.
Mit kleinen Schritten näherte sie sich Giorgias Bett. Sie hatte gehofft,
die einzige Besucherin zu sein. Das Einzelzimmer war zwar klein, dennoch
war um das Bett ausreichend Platz für mehrere Leute.

Giorgia schien zu schlafen. Valentina zog sich einen Stuhl direkt ans
Bett. Giorgias Arme lagen gerade ausgestreckt über der Decke neben ihr.
Trotz reichlich abbekommener Sonne während des Sommers lag ihr Gesicht
blass auf dem Kissen. Es wirkte wie Porzellan. Ihre dunklen, langen
Wimpern hoben sich scharf von der Haut ab. Valentina musste sich
zwingen, hinzusehen. Giorgia, dieser Ausbund an Vitalität, braungebrannt
selbst in den Wintermonaten, lag als Häuflein Elend vor ihr. Und wieder
kamen ein beklemmend trauriges Gefühl, Mitleid und der Gedanke hoch,
dass das alles doch gar nicht passiert sein konnte.
Plötzlich fing Giorgias Gesicht sachte an zu zucken. Valentina konnte
den Blick nicht mehr abwenden. Sie wartete und wartete und konnte
zusehen, wie sich Giorgias Augen zu schmalen Streifen öffneten. Es sah
beschwerlich aus.
Valentina rückte näher ans Bett, bis der Stuhl am Matratzenrand anstand.
Mit leiser Bewegung legte sie ihre Finger sanft auf Giorgias Hand.

«Giorgia?»

«Oh.»

Ihre halboffenen Augen lächelten.

«Wie schön.»

«Wie geht es dir?»

«Hm, wenn du, wenn du,» Giorgia holte Luft, «wenn du da, da bist, gut,»
damit drehte sie langsam ihre Hand, umfing Valentinas Finger, und nun
lächelte auch ihr Mund.

Für einen Moment sahen sie sich an. Giorgia schloss immer wieder die
Augen. Ihre Lider wirkten schwer.

«Seit wann bist du wach?»

Giorgia stöhnte.

«Weiss nicht.»

«Ich bin so froh, dass du wieder da bist.»

Valentina strich ihr eine Haarsträhne aus der Stirn.

«Ja.»

«Hast du Schmerzen?»

Giorgia drehte den Kopf zaghaft hin und her.

«Nein, die geben mir was, bin nur unendlich, unendlich müde.»

Während sie von freundlichen Ärzten, Krankenpflegerinnen- und Pflegern,
von ihren traurigen Eltern und treuen Freundinnen und Freunden erzählte,
brauchte sie etliche Pausen. Die wurden länger und länger. Valentina
hörte ihr still zu.
Mit einem leisen \emph{«ruh dich aus»} legte sich Valentina den
Zeigefinger vor die Lippen.

«Darf ich?»

Sie deutete Giorgia an, sich etwas an den Bettrand zu schieben, streifte
ihre Flipflops ab und legte sich vorsichtig neben sie aufs Bett.
Giorgia suchte nach Valentinas Hand. Wie zwei unbewegliche Puppen lagen
sie schliesslich nebeneinander.

«Mhm,» Valentina streichelte ihren Arm, «weisst du denn, wie das
passieren konnte?»

«Nein, ich erinnere mich fast an, fast an nichts.»

Jetzt starrte Giorgia an die Decke, schloss immer wieder die Augen und
rang nach Luft.

«Pillen, so was mit Pillen war da. Aber, aber, weiss nicht, ich, ich
erinnere mich nicht, ich hab doch nichts geschluckt.»

«Ich habe davon gehört.»

«Und Basti, Basti war da, da hinten bei mir, aber anders als, als
sonst.»

Valentinas Lächeln fror augenblicklich ein.

«Wie, er war da, wo denn?»

«Am Felsen, hinten am Felsen.»

«Da, wo du runtergefallen bist?»

«Hm, ja, ich glaube, er war da, aber,» Giorgia verzerrte das Gesicht,
«ich weiss nicht wann, alles, alles ist so,» das Sprechen schien sie zu
quälen, sie schnappte erneut nach Luft, «vielleicht auch nicht, oder
doch, Val, meine Er–, meine Erinnerung kommt nur langsam - »

Valentina setzte sich auf, legte Giorgia die flache Hand an die Wange
und stützte deren zur Seite gekippten Kopf ab.

«Alles gut, Giorgia, nur langsam, ruh dich aus.»

Giorgia sank tiefer in die Kissen. Sie schloss die Augen, atmete ein
und lange aus. Valentina schaute ihr dabei zu. In ihrem Kopf jagten sich die Gedanken. Was hatte Giorgia
gerade von Bastiano erzählt? Es musste etwas dran sein. Bastiano hatte
also am Unfallort tatsächlich Spuren hinterlassen?
Und Louis?

«Giorgia?»

Wieder öffnete sie schwerfällig die Augen.

«Hm?»

Valentina streichelte ihr übers Haar.

«War sonst noch jemand da?»

«Häh?»

Jetzt war Valentina, als würde sich Giorgia für einen Moment aus den
Kissen heben.
Sie schüttelte den Kopf energisch.

«Nein, niemand, vielleicht, vielleicht auch Basti nicht, oder, oder
doch. Weiss nicht genau, Val. Glaube, habe mich, der Ast, zu schwach,
der war zu schwach.»

Valentina war verwirrt. Sollte sie direkt nach Louis fragen? 
Oder würde sie Giorgia unnötig durcheinanderbringen? Sie liess es, legte sich
wieder flach hin und hoffte, ihre Freundin spreche weiter. Die machte
den Eindruck, sich wieder zu entfernen. Sie atmete regelmässiger.
Dann schlief sie langsam ein. Sie sah genau gleich aus, wie bei
Valentinas Ankunft.
Valentina erhob sich sacht und verliess das Zimmer. Kaum hatte sie
die Tür geschlossen, kam ihr ein Krankenpfleger entgegen. Er zielte auf
Giorgias Zimmertür.

«Buongiorno, Signora. Sie besuchten gerade Giorgia Sala?»

«Ja, ich bin eine Freundin. Sie ist gerade eingeschlafen.»

«Schön, haben Sie sie besucht. Was haben Sie für einen Eindruck von
ihr?»

«Es, es war wunderschön. Sie sprach mit mir, langsam zwar, aber ich verstand sie. »

«Das ist gut. Sie spricht erst seit gestern, wenig zwar und langsam, wie Sie sagen, aber bewusst.
Sie ist zurück.»


